\documentclass[10pt,a4paper]{amsart}
\usepackage[margin=19mm]{geometry}
\usepackage{amsmath,amssymb,mathtools}
\usepackage[scr=rsfs,cal=euler]{mathalfa}
\usepackage{xfrac}
\usepackage[unicode,hidelinks,bookmarks=false]{hyperref}
\newcommand{\ie}{i.e.\ }
\newcommand{\cf}{cf.\ }
\newcommand{\resp}{respectively\ }
\newcommand{\Subsection}{\subsection}
\providecommand{\nobreakdash}{\nobreak\hbox{-}}
\setlength{\emergencystretch}{2em}
\multlinegap0pt
\usepackage{bm}
\usepackage{amssymb}
\usepackage{breqn}
\usepackage{booktabs}
\usepackage{mathtools}
\usepackage{tikz}
\usepackage[all]{xy}
\usepackage{enumerate}
\newtheorem{lemma}{Lemma}
\newtheorem{proposition}{Proposition}
\newcommand{\C}{\mathbb{C}}
\renewcommand{\H}{\mathbf{h}}
\newcommand{\Ha}{\mathfrak{H}_{N}}
\newcommand{\Z}{\mathbb{Z}}
\title{Representations of Hamiltonian Lie algebras}
\author{Vyacheslav Futorny, Santanu Tantubay}
\date{Source: arXiv:2503.00749v3 (CC BY 4.0)}
\begin{document}
\maketitle

\noindent\emph{Source and licence.} Representations of Hamiltonian Lie algebras. Version arXiv:2503.00749v3 (CC BY 4.0), distributed by arXiv under the Creative Commons Attribution 4.0 International licence, http://creativecommons.org/licenses/by/4.0/, as recorded in the arXiv licence metadata for this exact version and shown on the abstract page https://arxiv.org/abs/2503.00749v3. Full licence notice: \texttt{licenses/arxiv-2503.00749-CC-BY-4.0.txt}.\par\smallskip

\noindent\textit{Editorial scope.} Counted unit: the article's proposition prop1 (source0.tex lines 286--293). The article's complete original proof of it is delivered with it (source0.tex lines 294--305, one \texttt{proof} environment of 12 lines). No mathematics is written, repaired or re-derived: the statement and the proof are the article's own lines, in the article's own notation, and no retained line is substituted or re-flowed.  Retained uncounted as the source-local context the proof needs: lines 180--182.  Nothing was omitted from any retained span, so the package carries no omission record.  No retained line contains a citation, so the wrapper carries no bibliography and no bibliography entry.  No retained line contains a figure environment, an image-inclusion command or drawing code, so no image is delivered and the package declares no asset dependency.  Editorial formatting only, in addition: an AMS \texttt{amsart} preamble that restates the article's own declarations and macros used by the retained lines, namely the environments \texttt{lemma}, \texttt{proposition} on separate counters, together with the standard packages those lines need.  The retained environments print in this document as Lemma 1, Proposition 1 in order of appearance, whereas the article prints the same items with the numbering of its own full document.  The complete native source of the article as distributed in the arXiv e-print for this exact version is bundled unchanged as \texttt{sources/174096/source0.tex}.
\begin{lemma}\label{la}
   For each $a\in \Z_+^N$ consider $c_a\in End(V)$ and assume that only finitely many $c_a$ are nonzero. Define $g(m)=\sum_{a\in \Z_+^N}c_am^a$,  $m\in \Z^N$. Let $M$ be a subspace of $V$.  If $g(m)v\in M$ for some $v\in V$ and all $m\in \Z^N$, then $c_av\in M$ for all $a\in \Z_+^N$. 
\end{lemma}

\begin{proposition}\label{prop1}
   The $\Ha$-module $F^{\alpha,\beta}(V(\delta_0))=\C\otimes A_N$ is irreducible if and only if  $\alpha\notin \Z^N$.
   
   
    
%    Let $\alpha\in \C^N\setminus \Z^N$, then $F^\alpha(V(\delta_0))$ is irreducible module over $\Ha$. 
    
\end{proposition}

\begin{proof}
Suppose that $\alpha\in \Z^N$. Then $F^{\alpha,\beta}(V(\delta_0))$ is reducible  since $\C(1\otimes t^{-\alpha})$ is a nonzero proper submodule of $F^\alpha(V(\delta_0))$. It is easy to see that the quotient module is irreducible.
 Let $\alpha\in \C^N\setminus \Z^N$. We will show that $F^{\alpha,\beta}(V(\delta_0))$ is irreducible. 
    Let $P$ be nonzero submodule of $F^{\alpha,\beta}(V(\delta_0))$. 
    Now it being a weight module, there exists $k\in \Z^N$ such that $1\otimes t^k\in P$. Suppose $r\in \Z^N$ such that $\bm{(}\bar{r}, k+\alpha\bm{)}\neq 0$, then $\H_r(1\otimes t^k)=\bm{(}\bar{r}, k+\alpha\bm{)}(1\otimes t^{k+r})\in P$ i.e. $1\otimes t^{k+r}\in P$. Now we choose a fixed $r\in \Z^N$ such that $\bm{(}\bar{r}, k+\alpha\bm{)}=0$. 
    
    Since $\alpha\notin \Z^N $, we will have $(k+\alpha)\neq 0$, so there exists $i\in \{1,\dots, 2n \}$ such that $(k_i+\alpha_i)\neq 0$. Now $\H_{r-s}\H_s(1\otimes t^k)\in P$ for all $r,s\in \Z^N$. This means $g_2(s)\cdot 1\in P_{k+r}$ for all $r\in \Z^N$. Now we see that
     \[ 
      g_2(s)\cdot 1=\bm{(}\bar{r}-\bar{s},k+s+\alpha\bm{)}\bm{(}\bar{s},k+\alpha\bm{)}1=(\bm{(}\bar{r},s\bm{)}-\bm{(}\bar{s},k+\alpha\bm{)})\bm{(}\bar{s},k+\alpha\bm{)}1\]
      \[=-\bm{(}\bar{s},r+k+\alpha\bm{)}\bm{(}\bar{s},k+\alpha\bm{)}1,\]
since we chose $r\in \Z^N$ such that $\bm{(}\bar{r},k+\alpha\bm{)}=0$. Now, as $\alpha\notin \Z^N$, there exist $1\leq i,j\leq N$ such that $(r_i+k_i+\alpha_i)(k_j+\alpha_j)\neq 0$. So we will have a second degree monomial of $g_2(s)$, whose coefficient is $-(r_i+k_i+\alpha_i)(k_j+\alpha_j)$. Applying Lemma~\ref{la}, we can say that $-(r_i+k_i+\alpha_i)(k_j+\alpha_j) 1\in P_{k+r}$. Combining both  cases, we conclude that $1\in P_{k+r}$ for all $r\in \Z^N$. This implies that $P=F^{\alpha,\beta}(V(\delta_0))$ and proves the irreducibility.
\end{proof}

\end{document}
